\pdfoutput=1
\documentclass[11pt]{amsart}
\usepackage[margin=1.2in]{geometry}
\usepackage{amsmath,amssymb,amsthm}
\usepackage[hidelinks]{hyperref}
\hypersetup{pdftitle={The universality of words x^r y^s in symmetric groups},
  pdfauthor={Avi Kondareddy}}
\numberwithin{equation}{section}
\raggedbottom

\newtheorem{theorem}{Theorem}[section]
\newtheorem{lemma}[theorem]{Lemma}
\newtheorem{corollary}[theorem]{Corollary}
\newtheorem{proposition}[theorem]{Proposition}
\theoremstyle{remark}
\newtheorem{remark}[theorem]{Remark}
\newtheorem*{mainthm}{Main Theorem}

\newcommand{\Sym}{\operatorname{Sym}}
\newcommand{\supp}{\operatorname{supp}}
\newcommand{\sgn}{\operatorname{sgn}}
\newcommand{\md}{\mathrm{m}}   % moved points
\newcommand{\cy}{\mathrm{c}}   % nontrivial cycles

\title{The universality of words $x^ry^s$ in symmetric groups}
\author{Avi Kondareddy}
\date{August 25, 2026}
\subjclass[2020]{Primary 20B30; Secondary 20D60, 11P32, 05A05}
\keywords{word maps, universal words, symmetric groups, products of cycles,
three-primes theorem}

\begin{document}

\begin{abstract}
We answer Problem~10.32 of the Kourovka Notebook, posed by Yu.~I.~Merzlyakov in
1986: for all nonzero integers $r,s$ at least one of which is odd, the word
$x^ry^s$ is universal on the symmetric group $S_n$ --- every permutation is a
value --- whenever $n\ge\max\{N_0,\,C\log m\}$, where $m=m(r,s)$ is the largest
squarefree divisor of $rs$ and $C,N_0$ are absolute constants. This is the
symmetric-group analogue of the alternating-group theorem of Brenner, Evans and
Silberger, and it replaces Droste's bound $n\ge\max\{6,\,4m-4\}$, linear in $m$
and the best known since 1986. Even targets, and odd targets moving at least
$n/2$ points, are factored as products of two cycles of explicitly controlled
prime-related lengths, via the Herzog--Kaplan--Lev two-cycle criterion and
Chebyshev-type prime counting; this part of the argument is explicit and fully
effective. For odd targets moving fewer than $n/2$ points we show that, for
adversarial $r,s$, any such two-cycle factorization would encode prime
constellations of twin-prime type. These targets are instead factored, inside a
suitable subset of $\{1,\dots,n\}$, as a permutation of cycle type
$(q_1,q_2,q_3,1)$ times a full cycle of even length, via Boccara's criterion for
products with a full cycle; the primes $q_1,q_2,q_3$, none dividing $rs$, come
from Vinogradov's three-primes theorem relative to the excluded set of primes of
$rs$ --- the one ineffective ingredient. A complementary construction respecting
the parity hypothesis ($r=\operatorname{lcm}(1,\dots,n)$, $s$ its odd part)
forces every value of the word to have all cycle lengths powers of two; the
constant therefore cannot be taken to be $1$, and the logarithmic shape is
optimal for $S_n$, as it is for $A_n$.
\end{abstract}

\maketitle

\section{Introduction}\label{sec:intro}

For which degrees $n$ is every permutation in $S_n$ a product of an $r$-th power
and an $s$-th power? In the language of word maps: a word $w$ in the free group
on $x,y$ is \emph{universal} on a group $G$ if the map $(x,y)\mapsto w(x,y)$
sends $G\times G$ onto $G$, and the question asks when the two-letter power word
$x^ry^s$, $rs\neq0$, is universal on $S_n$. Surjectivity questions for word maps
are the noncommutative Waring problems of group theory: Ore's conjecture ---
every element of a finite non-abelian simple group is a commutator, that is,
the word $[x,y]$ is universal --- stood for nearly six decades before its
resolution by Liebeck, O'Brien, Shalev and Tiep \cite{LOST}. For an arbitrary
nontrivial word $w$, the theorems of Larsen--Shalev and Larsen--Shalev--Tiep
\cite{LaSh,LST} show that the values of $w$ cover every sufficiently large such
group in two steps: $w(G)^2=G$. For the power words one can ask for more --- genuine universality,
in one step --- and the classical literature on products of conjugacy classes
of $S_n$ makes the question concrete.

What no general theorem gives is uniformity in the word. For $w=x^ry^s$ the
threshold $N(w)$ of \cite{LaSh} carries no dependence on $(r,s)$ that one can
read off, and the exponents can conspire against any fixed degree
(\S\ref{sec:lower} exhibits the conspiracy). The right question, posed forty
years ago, is quantitative: how must the degree grow with the arithmetic of
$(r,s)$ for universality to set in? Write $m=m(r,s)$ for the largest squarefree
divisor of $rs$ (so $m$ is the product of the distinct primes dividing $rs$),
with the convention $m=1$ when $r,s\in\{-1,1\}$; the answer, it turns out, is
governed by $m$ alone, and is logarithmic in it.

Two papers in the same 1986 issue of the \emph{Proceedings of the AMS} frame the
problem. Brenner, Evans and Silberger \cite{BES} proved that for every constant
$c>8/5$ there is an $N_0(c)$ such that $x^ry^s$ is universal on $A_n$ whenever
$n\ge\max\{N_0(c),\,c\log m\}$; they showed moreover that the constant cannot be
taken to be $1$, so the logarithmic shape is optimal for $A_n$. Droste \cite{Droste}
proved, among other things, that when at least one of $r,s$ is odd the word
$x^ry^s$ is universal on $S_n$ for $n\ge\max\{6,\,4m-4\}$ --- a bound linear in
$m$, hence exponentially weaker in $\log m$. Problem~10.32 of the Kourovka
Notebook, contributed by Yu.~I.~Merzlyakov \cite{KN} to its tenth issue (1986)
--- the same year as the two \emph{Proceedings} papers --- asks:

\begin{quote}
\emph{Find an analogous \textup{[}logarithmic\textup{]} bound for the degree of the
symmetric group $S_n$ under hypothesis that at least one of the $r,s$ is odd:
up to now, here only a more crude bound $n\ge\max\{6,4m(r,s)-4\}$ is known
\textup{(}M.~Droste, Proc.\ Amer.\ Math.\ Soc., 96, no.~1 (1986),
18--22\textup{)}.}
\end{quote}

The parity hypothesis is necessary: if $r$ and $s$ are both even then every value
$x^ry^s$ is an even permutation, and the word is not universal on $S_n$ for any
$n\ge 2$. This paper answers the problem in both directions: the logarithmic
bound holds, and the logarithmic shape is optimal.

\begin{mainthm}[answer to Kourovka Notebook Problem 10.32]
There exist absolute constants $C$ and $N_0$ with the following property. For all
nonzero integers $r,s$ at least one of which is odd, the word $x^ry^s$ is universal
on the symmetric group $S_n$ whenever
\[
n\;\ge\;\max\{N_0,\;C\log m(r,s)\}.
\]
\end{mainthm}

The logarithmic threshold is the natural one. The word sees only the primes
dividing $rs$: a cycle of length coprime to $r$ is an $r$-th power
(Lemma~\ref{lem:root}), so what the proofs must manufacture are cycle lengths up
to $n$ avoiding a prescribed set of primes. That set has total logarithmic mass
$\log m=\sum_{p\mid rs}\log p$, while by Chebyshev the primes up to $n$ have
mass about $n$; so $n\asymp\log m$ is exactly the point at which the primes of
$rs$ can no longer swallow every prime below $n$. The lower bound of
\S\ref{sec:lower} shows the heuristic is sharp on the other side: taking
$r=\operatorname{lcm}(1,\dots,n)$ (and $s$ its odd part) makes every prime up
to $n$ divide $rs$ at cost $\log m=\theta(n)$, and $\theta(n)<n$ infinitely
often --- the arithmetic of $rs$ really can cover every degree up to $n$.

The constant $C$ produced by the proof is explicit but far from optimal; the
threshold $N_0$ is ineffective, for the single reason that the proof invokes
Vinogradov's three-primes theorem (see Remark~\ref{rem:effective}). For even
targets, and for odd targets moving at least $n/2$ points, the proof is fully
explicit and effective and gives $n\ge\max\{N_0',\,4.65\log m\}$ with an explicit
$N_0'$; it is only the remaining case --- odd targets of small support --- that
requires analytic input. For the complementary lower bound, an adaptation of the
Brenner--Evans--Silberger examples that respects the parity hypothesis
(\S\ref{sec:lower}) shows the constant cannot be taken to be $1$, so the
logarithmic shape is optimal for $S_n$ just as for $A_n$.

\subsection*{Why the symmetric group resisted}
The argument of Brenner, Evans and Silberger runs through elements of order
coprime to
$m$: if every $z\in A_n$ is a product of two such elements, one wins by taking
$r$-th and $s$-th roots inside the cyclic groups they generate. On $S_n$ this
mechanism breaks at once for \emph{odd} targets: an odd permutation has at least
one even-length cycle, so its order is even, and when $2\mid m$ no odd permutation
of order coprime to $m$ exists at all. The two factors must therefore be treated
asymmetrically --- the first need only be an $r$-th power, the second an $s$-th
power --- and this is where the hypothesis enters: with $s$ odd, even-length cycles
of length coprime to $s$ exist, and they carry the odd sign.

A second, subtler obstruction governs the proof: the first dictated the
asymmetry of the factors; this one dictates their shape. Products of \emph{two
single cycles} are completely described by a criterion of Herzog, Kaplan and Lev
(Theorem~\ref{thm:HKL} below): $z$ is a product of an $l_1$-cycle and an
$l_2$-cycle essentially if and only if $l_1+l_2$ is at least $\md(z)+\cy(z)$ and of
the same parity, and $|l_1-l_2|\le \md(z)-\cy(z)$, where $\md(z)$ is the number of
moved points and $\cy(z)$ the number of nontrivial cycles. For an odd target the
two cycle lengths must have opposite parities, so $|l_1-l_2|$ is odd and in
particular nonzero; when $\delta(z):=\md(z)-\cy(z)$ is small --- think of a single
transposition times a $3$-cycle --- the difference constraint forces the odd length
$u$ and even length $v$ into an interval of bounded width, and $u$ must avoid the
prime divisors of $r$, $v$ those of $s$. For adversarial $r,s$ this demands, in
effect, a pair of primes a bounded distance apart in prescribed positions --- a
constellation of twin-prime or binary-Goldbach strength, beyond reach. (Bounded
gaps between primes are known after Zhang and Maynard, but not in prescribed
windows, and the pair needed here involves the dilate $2^aq$ rather than a
shift.)
\emph{No choice of two single cycles can be proved to work for all such targets
short of solving such a constellation problem}; this is made precise in
Proposition~\ref{prop:obstruction}. The problem thus sits on the classical
watershed of additive prime number theory: two cycle lengths give a binary
constraint, of twin-prime strength, while a third summand would move the problem
to ternary Goldbach, where Vinogradov's theorem waits.

Our escape buys that third summand with a different factor shape. For odd $z$ of
small support we write $z=AB$ inside an $N$-point subset, with $N=2\tilde q$ for
a prime $\tilde q\nmid s$. The factor $B$ is a single $N$-cycle, hence an
$s$-th power; the factor $A$, an $r$-th power, has cycle type $(q_1,q_2,q_3,1)$
with $q_1,q_2,q_3$ primes summing to $N-1$, each avoiding the primes of $rs$.
The existence of such a factorization is governed by a beautiful criterion of
Boccara (Theorem~\ref{thm:Boccara}): a bound
$\ell(\lambda)+\ell(\mu)\le N+1$ --- a genus bound, via the branched-cover
correspondence noted in \S\ref{sec:prelim} --- together with a parity condition. The genus bound
demands exactly $\delta(z)\ge 3$ --- and every odd non-transposition satisfies it
--- while the parity condition is automatic for odd $z$. The three primes are
produced by Vinogradov's theorem \emph{relative to an excluded set}: the number of
representations $N-1=q_1+q_2+q_3$ is $\gg N^2/\log^3N$, while a sieve bound,
averaged over the $\asymp n/\log n$ admissible choices of $\tilde q$, shows that
representations meeting the primes of $rs$ are strictly fewer for at least one
choice (indeed for most). The order of quantifiers is essential: $r,s$ are fixed
before $\tilde q$ is chosen, so the excluded set cannot correlate with $N$.

Three elementary observations make the case division clean. First, every
nontrivial cycle
moves at least two points, so $\cy(z)\le \md(z)/2$ and hence
$\delta(z)\ge \md(z)/2$: \emph{small $\delta$ forces small support}, and the
troublesome targets of the previous paragraph automatically leave at least half of
$\{1,\dots,n\}$ free. Second, $\sgn(z)=(-1)^{\delta(z)}$, so $\delta$ is odd
exactly for odd $z$, which is what makes the parity conditions in both criteria
fall into place. Third, an odd $z$ with $\delta(z)=1$ is a bare transposition
$\tau$, and then $y=\tau$, $x=1$ works outright since $\tau^s=\tau$ for odd $s$.

\subsection*{Structure of the proof}
After reductions (\S\ref{sec:red}) we may take $s$ odd and both exponents at least
$2$ in absolute value. The theorem is then proved by three constructions:
\begin{enumerate}
\item[(I)] \emph{$z$ even, $z\neq 1$} (\S\ref{sec:even}): $z=(\text{$q$-cycle})
\cdot(\text{$q$-cycle})$ for a single prime $q\nmid rs$ in $(3n/4,\,n]$. On the
diagonal $l_1=l_2$ the Herzog--Kaplan--Lev difference constraint is vacuous, and
Chebyshev-type counting supplies the prime as soon as $n\ge4.5\log m$.
\item[(II)] \emph{$z$ odd, $\md(z)\ge n/2$} (\S\ref{sec:dense}): $z=(\text{$u$-cycle})
\cdot(\text{$v$-cycle})$ with $u$ a prime not dividing $r$ in $(3n/4,\,n]$ and
$v=2^aq'$ for a prime $q'\nmid s$ in a dyadic rescaling
$(3n/2^{a+2},\,n/2^a]$ of the same window, $1\le a\le5$. Here
$|u-v|<n/4\le\delta(z)$ and $u+v>3n/2\ge \md(z)+\cy(z)$, and the counting needs
$n\ge4.65\log m$.
\item[(III)] \emph{$z$ odd, $\md(z)<n/2$, $\delta(z)\ge3$} (\S\ref{sec:sparse}):
the Boccara factorization described above, on an $N$-point set with
$N=2\tilde q\in(n/2,\,n]$.
\end{enumerate}
Section~\ref{sec:proof} assembles the theorem, \S\ref{sec:lower} proves the
complementary lower bound, and \S\ref{sec:remarks} closes with a synthesis and
remarks on the shapes of the factors, effectivity, and constants.

\subsection*{Relation to the literature}
As discussed above, the general word-map theorems \cite{LaSh,LST} are
qualitative in the parameters of the word: for $w=x^ry^s$ they give no
uniformity of $N(w)$ in $(r,s)$, and in particular no bound of the shape
$n\gg\log m(r,s)$, which is the point of Problem~10.32. The complementary
literature on explicit factorization --- products of two prescribed conjugacy
classes (Bertram \cite{Bertram}; Boccara \cite{Boccara82}; Herzog, Kaplan and
Lev \cite{HKL}; the survey material in \cite{Petronio}) --- provides exactly the
combinatorial control needed here, and this paper joins those criteria to
classical analytic number theory at a single point. The content lies in locating
that point --- Proposition~\ref{prop:obstruction} shows the analytic input is
unavoidable for any two-single-cycle approach, not merely convenient --- and in
choosing a factor shape for which the input needed is ternary rather than
binary.

\section{Preliminaries}\label{sec:prelim}

\subsection{Notation}
Throughout, $n$ denotes the degree of the symmetric group $S_n$ and $z\in S_n$ a
target permutation. We write $\supp(z)$ for the set of moved points,
$\md(z)=|\supp(z)|$, $\cy(z)$ for the number of cycles of length at least $2$, and
\[
\delta(z)\;=\;\md(z)-\cy(z).
\]
A \emph{cycle type} on an $N$-point set $\Omega$ is a partition $\lambda\vdash N$
in which fixed points are recorded as parts equal to $1$; we write $\ell(\lambda)$
for the total number of parts, fixed points included, and $C_\lambda$ for the
conjugacy class of $\Sym(\Omega)$ with that type. For $r,s$ we set
$m=m(r,s)$ as in the introduction and
\[
L\;=\;\log m\;=\;\sum_{p\mid rs}\log p ,
\]
sums and products over $p$ always ranging over primes. All logarithms are natural.

\begin{lemma}\label{lem:elem}
For every $z\in S_n$:
\begin{enumerate}
\item $\sgn(z)=(-1)^{\delta(z)}$; in particular $z$ is odd if and only if
$\delta(z)$ is odd.
\item $\cy(z)\le \md(z)/2$, hence $\delta(z)\ge \md(z)/2\ge \cy(z)$.
\item If $z$ is odd and $\delta(z)=1$ then $z$ is a transposition.
\end{enumerate}
\end{lemma}

\begin{proof}
An $\ell$-cycle has sign $(-1)^{\ell-1}$; multiplying over the nontrivial cycles
of $z$, whose lengths sum to $\md(z)$ and whose number is $\cy(z)$, gives (1).
Each nontrivial cycle moves at least $2$ points, which is (2). For (3),
$\delta(z)=1$ and (2) give $\md(z)\le 2$, so $z$ moves exactly two points.
\end{proof}

\subsection{Roots}
The next lemma is the standard fact about powers of cycles (see, e.g.,
\cite[\S4.8]{Wilf}: a cycle of length $\ell$ raised to the $t$-th power is a
product of $\gcd(\ell,t)$ cycles of length $\ell/\gcd(\ell,t)$).

\begin{lemma}\label{lem:root}
Let $\sigma\in S_n$ be a product of pairwise disjoint cycles of lengths
$\ell_1,\dots,\ell_k$, and let $t$ be a nonzero integer with $\gcd(\ell_i,t)=1$
for all $i$. Then there is $x\in S_n$ with $\supp(x)=\supp(\sigma)$ and
$x^t=\sigma$.
\end{lemma}

\begin{proof}
It suffices to treat a single $\ell$-cycle $\gamma$ and multiply the resulting
roots, which are disjoint. Choose $a$ with $at\equiv 1\pmod{\ell}$; then
$(\gamma^a)^t=\gamma^{at}=\gamma$ since $\gamma^\ell=1$, and $\gamma^a$ is again an
$\ell$-cycle on the same support because $\gcd(a,\ell)=1$.
\end{proof}

\subsection{Two factorization criteria}
The first criterion characterizes products of two single cycles. It was proved by
Herzog, Kaplan and Lev; the case of two cycles of given lengths goes back to
Boccara \cite{Boccara80}, and the two-$l$-cycle case to Bertram \cite{Bertram}.

\begin{theorem}[Herzog--Kaplan--Lev {\cite[Theorem~7]{HKL}}]\label{thm:HKL}
Let $\sigma\in S_n$, $\sigma\neq 1$, and let $n\ge l_1\ge l_2\ge 2$ be integers.
Then $\sigma$ is a product of an $l_1$-cycle and an $l_2$-cycle in $S_n$ if and
only if one of the following holds:
\begin{enumerate}
\item[(a)] $\sigma$ has exactly two nontrivial cycles, and their lengths are $l_1$
and $l_2$; or
\item[(b)] $l_1+l_2\ge \md(\sigma)+\cy(\sigma)$, \quad
$l_1+l_2\equiv \md(\sigma)+\cy(\sigma) \pmod 2$, \quad and \quad
$l_1-l_2\le \md(\sigma)-\cy(\sigma)$.
\end{enumerate}
\end{theorem}

We use the sufficiency of (b) throughout; the necessity direction enters only in
the obstruction Proposition~\ref{prop:obstruction}. Note that if $\sigma=AB$ then
also $\sigma=B\,(B^{-1}AB)$, so the two lengths may be realized in either order; we
use this freedom silently to put the $r$-th-power factor first.

The second criterion characterizes products of an arbitrary class with the class
of full cycles. It is due to Boccara \cite{Boccara82}; it is also equivalent, via
the Hurwitz monodromy correspondence, to the realizability of branched covers of
the sphere with three branch points one of which is totally ramified, proved by
Thom for the sphere and by Edmonds, Kulkarni and Stong in general \cite{Thom,EKS}
(see \cite[Theorem~2.1]{Petronio} for a modern statement).

\begin{theorem}[Boccara {\cite{Boccara82}}]\label{thm:Boccara}
Let $\Omega$ be a set of size $N\ge 2$, let $\lambda\vdash N$ be a cycle type
(fixed points counted as parts $1$), and let $z\in\Sym(\Omega)$ have cycle type
$\mu\vdash N$. Then $z=AB$ for some $A\in C_\lambda$ and some $N$-cycle $B$ on
$\Omega$ if and only if
\[
\ell(\lambda)+\ell(\mu)\;\le\;N+1
\qquad\text{and}\qquad
\ell(\lambda)+\ell(\mu)\;\equiv\;N+1\;\pmod 2 .
\]
\end{theorem}

\begin{remark}\label{rem:machine}
Both criteria have been verified mechanically by exhaustive computation:
Theorem~\ref{thm:HKL} over all of $S_7$ and $S_8$ (all $1618$ triples (cycle
type, $l_1$, $l_2$) with $2\le l_1,l_2\le n$, the identity type included as a
sanity case), and Theorem~\ref{thm:Boccara} for all $N\le 8$ (all ordered pairs
$(\lambda,\mu)$ of types, both directions of the equivalence), with no
discrepancies. We mention this because both statements are used here in
load-bearing ways and details of \cite{Boccara82} circulate imprecisely --- its
volume number, for instance, appears incorrectly in more than one recent
bibliography.
\end{remark}

\subsection{Prime counting}
Let $\theta(x)=\sum_{p\le x}\log p$ and $\pi(x)$ the usual counting functions. We
use the classical explicit bounds of Rosser and Schoenfeld: \eqref{eq:RSup} is
\cite[Theorem~9, eq.~(3.32)]{RS}, \eqref{eq:piup} and \eqref{eq:pilow} are
\cite[eqs.~(3.6) and (3.5)]{RS} respectively, and \eqref{eq:RSlow} is from their
later paper
\cite[Corollary to Theorem~6]{RS75}:
\begin{align}
\theta(x)&\;<\;1.01624\,x  &&\text{for all } x>0, \label{eq:RSup}\\
\theta(x)&\;>\;0.985\,x    &&\text{for } x> 11927, \label{eq:RSlow}\\
\pi(x)&\;<\;1.25506\,\frac{x}{\log x} &&\text{for } x>1, \label{eq:piup}\\
\pi(x)&\;>\;\frac{x}{\log x} &&\text{for } x\ge 17. \label{eq:pilow}
\end{align}

\begin{lemma}[Prime windows]\label{lem:window}
Let $0<\alpha<\beta\le 1$ and put $K=0.985\,\beta-1.01624\,\alpha$, assumed
positive. Let $Q$ be any set of primes with $\sum_{q\in Q}\log q\le L$. If
$n> 11927/\beta$ and $Kn>L$, then the interval $(\alpha n,\,\beta n]$ contains a
prime not belonging to $Q$.
\end{lemma}

\begin{proof}
By \eqref{eq:RSup} and \eqref{eq:RSlow},
$\sum_{\alpha n<p\le\beta n}\log p=\theta(\beta n)-\theta(\alpha n)
> 0.985\beta n-1.01624\alpha n=Kn>L$. The primes of $Q$ in the window contribute
at most $\sum_{q\in Q}\log q\le L$ to the left side, so some prime of the window
lies outside $Q$.
\end{proof}

We record the windows we use, with $K$ rounded down:
\begin{equation}\label{eq:windows}
\begin{gathered}
(\tfrac34 n,\,n]:\ K=0.2228;\\
\text{rescaled, } (3n/2^{a+2},\,n/2^a]:\ K=0.2228/2^a
\quad(a\ge1,\ n>11927\cdot2^a).
\end{gathered}
\end{equation}
In particular the base window contains a prime avoiding any prescribed set of
primes of log-weight at most $L$ provided $n\ge 23856$ and $n\ge 4.5\,L$; the
rescaled windows are used, five at a time, in \S\ref{sec:dense}.

\begin{lemma}\label{lem:omega}
For every integer $m\ge 2$ with $L=\log m$,
\[
\omega(m)\;\le\;1.28\,\frac{L}{\log L}\;+\;1429 ,
\]
where $\omega(m)$ is the number of distinct prime divisors of $m$ (and the first
term is read as $0$ when $L\le 1$).
\end{lemma}

\begin{proof}
Let $\omega=\omega(m)$ and let $q_\omega$ denote the $\omega$-th prime. The $i$-th
smallest prime divisor of $m$ is at least the $i$-th prime, so
$L\ge\theta(q_\omega)$. If $q_\omega\le11927$ then
$\omega=\pi(q_\omega)\le\pi(11927)=1429$. Otherwise $q_\omega>11927$ and
\eqref{eq:RSlow} gives
$q_\omega\le L/0.985\le 1.0153\,L$, and \eqref{eq:piup} gives
$\omega=\pi(q_\omega)\le\pi(1.0153L)<1.25506\cdot
1.0153\,L/\log(1.0153L)<1.28\,L/\log L$.
\end{proof}

\subsection{Analytic input}
For a positive integer $k$ let
\begin{align*}
r_2(k)&=\#\{(p_1,p_2):\ p_i \text{ prime},\ p_1+p_2=k\},\\
r_3(k)&=\#\{(p_1,p_2,p_3):\ p_i \text{ prime},\ p_1+p_2+p_3=k\}
\end{align*}
count ordered representations as sums of primes. We use the following two
classical facts.

\begin{theorem}[Vinogradov \cite{Vino}; see {\cite[Theorem~3.4]{Vaughan}}]
\label{thm:vino}
There are constants $c_1>0$ and $M_0$ such that $r_3(M)\ge c_1 M^2/\log^3 M$ for
every odd $M\ge M_0$.
\end{theorem}

\begin{theorem}[Sieve bound; see {\cite[Theorem~3.11]{HR}}]\label{thm:sieve}
There is an absolute constant $C_2$ such that for every even $k\ge 4$,
\[
r_2(k)\;\le\;C_2\,S(k)\,\frac{k}{\log^2 k},
\qquad\text{where}\quad
S(k)\;=\;\prod_{p\mid k,\ p>2}\frac{p-1}{p-2} .
\]
For odd $k$, trivially $r_2(k)\le 2$ (one summand must be $2$).
\end{theorem}

We also use the Brun--Titchmarsh inequality in the form of Montgomery and Vaughan
\cite[Theorem~2]{MV}: for coprime $a,d$ and $x>d$,
\begin{equation}\label{eq:BT}
\pi(x;d,a)\;<\;\frac{2x}{\varphi(d)\log(x/d)} ,
\end{equation}
where $\pi(x;d,a)$ counts the primes $p\le x$ with $p\equiv a\pmod d$.

\begin{remark}
Theorem~\ref{thm:vino} is the one ineffective ingredient of this paper: $M_0$ is
not computable from the classical proof (Siegel--Walfisz). The constant $c_1$ is
an explicit number, and effective versions of the three-primes theorem exist
(Helfgott \cite{Helfgott}); see Remark~\ref{rem:effective}.
\end{remark}

\section{Reductions}\label{sec:red}

\begin{lemma}\label{lem:red}
Let $G$ be a group and $r,s$ nonzero integers.
\begin{enumerate}
\item $x^ry^s$ is universal on $G$ if and only if $x^sy^r$ is.
\item Universality of $x^ry^s$ on $G$ depends only on $(|r|,|s|)$.
\item If $|r|=1$ or $|s|=1$ then $x^ry^s$ is universal on every group.
\end{enumerate}
\end{lemma}

\begin{proof}
(1) Given $z\in G$, apply universality of $x^sy^r$ to $z^{-1}$: if
$z^{-1}=a^sb^r$ then $z=b^{-r}a^{-s}=(b^{-1})^r(a^{-1})^s$. (2) Substitute
$x\mapsto x^{-1}$, $y\mapsto y^{-1}$ as needed. (3) If $|s|=1$ take $x=1$,
$y=z^{\pm1}$; if $|r|=1$ take $y=1$, $x=z^{\pm 1}$.
\end{proof}

Since $m(r,s)=m(s,r)$ and $m$ depends only on $(|r|,|s|)$, Lemma~\ref{lem:red}
allows us to assume for the remainder of the paper:
\begin{equation}\label{eq:standing}
r\ge 2,\qquad s\ge 3 \text{ odd},
\end{equation}
the excluded cases being trivially universal. (If the odd exponent among $r,s$ is
the first one, swap by Lemma~\ref{lem:red}(1); then drop signs by (2); if either
exponent is $1$ use (3), which in particular covers $m=1$.)

Two families of targets are handled at once. If $z=1$, take $x=y=1$. If $z$ is a
transposition $\tau$, take $x=1$ and $y=\tau$: since $s$ is odd,
$y^s=\tau^s=\tau=z$. By Lemma~\ref{lem:elem}(3) these dispose of all odd $z$ with
$\delta(z)=1$.

\section{Even targets}\label{sec:even}

\begin{proposition}\label{prop:even}
Assume \eqref{eq:standing}, $n\ge 23856$ and $n\ge4.5\,L$. Then every even
$z\in S_n$, $z\neq 1$, satisfies $z=x^ry^s$ for some $x,y\in S_n$.
\end{proposition}

\begin{proof}
Under \eqref{eq:standing} the odd number $s\ge3$ contributes an odd prime to $m$,
so $m\ge 3$, $L\ge\log 3>0$, and
$Kn\ge 0.2228\cdot4.5\,L>L$ for the window $(\tfrac34n,n]$ of
\eqref{eq:windows}. By Lemma~\ref{lem:window}, applied to this window with $Q$
the set of primes dividing $rs$, there is a prime
\[
q\in(\tfrac34 n,\,n],\qquad q\nmid rs .
\]
We check condition (b) of Theorem~\ref{thm:HKL} for $l_1=l_2=q$. Since
$\cy(z)\le \md(z)/2$ (Lemma~\ref{lem:elem}) and $\md(z)\le n$,
\[
l_1+l_2=2q>\tfrac32n\ \ge\ \md(z)+\cy(z) .
\]
As $z$ is even, $\delta(z)=\md(z)-\cy(z)$ is even, so $\md(z)+\cy(z)$ is even as
well, matching $l_1+l_2=2q$. Finally $l_1-l_2=0\le \md(z)-\cy(z)$, the latter
being positive since $z\neq 1$. Hence $z=AB$ with $A,B$ both $q$-cycles. Since
$q\nmid rs$, Lemma~\ref{lem:root} provides $x$ with $x^r=A$ and $y$ with $y^s=B$,
and $x^ry^s=AB=z$.
\end{proof}

\section{Odd targets of large support}\label{sec:dense}

\begin{proposition}\label{prop:dense}
Assume \eqref{eq:standing}, $n\ge 381700$ and $n\ge4.65\,L$. Then every odd
$z\in S_n$ with $\md(z)\ge n/2$ satisfies $z=x^ry^s$ for some $x,y\in S_n$.
\end{proposition}

\begin{proof}
Write $\delta=\md(z)-\cy(z)$. By Lemma~\ref{lem:elem}(2),
$\delta\ge \md(z)/2\ge n/4$. As in Proposition~\ref{prop:even}, $L>0$, so
$0.2228\,n\ge0.2228\cdot4.65\,L>L$, and Lemma~\ref{lem:window} applied to the
window $(\tfrac34n,n]$ with $Q=\{p:p\mid r\}$ yields a prime
\[
u\in(\tfrac34n,\,n],\qquad u\nmid r
\]
(note $u$ is odd since $u>3$).

For the even factor we use the dyadic windows of \eqref{eq:windows}. Suppose
that for every $a\in\{1,\dots,5\}$, every prime in
$W_a=(3n/2^{a+2},\,n/2^a]$ divides $s$. The intervals $W_a$ are pairwise
disjoint (the top of $W_{a+1}$ is $n/2^{a+1}<3n/2^{a+2}$, the bottom of $W_a$),
and $n/2^a>11927$ for $a\le5$ since $n\ge381700>11927\cdot2^5$; so, applying
\eqref{eq:RSup} and \eqref{eq:RSlow} to each window and summing,
\[
L\;\ge\;\sum_{a=1}^{5}\bigl(\theta(n/2^a)-\theta(3n/2^{a+2})\bigr)
\;\ge\;\sum_{a=1}^{5}\frac{0.2228}{2^a}\,n\;=\;\frac{31}{32}\cdot0.2228\,n
\;>\;0.2158\,n ,
\]
contradicting $n\ge4.65\,L$ (as $4.65\cdot0.2158>1$). Hence there are $a\le5$
and a prime $q'\in W_a$ with $q'\nmid s$; put
\[
v\;=\;2^aq'\;\in\;(\tfrac34n,\,n] ,
\]
an even number with $\gcd(v,s)=1$ ($s$ is odd and $q'\nmid s$).

We check Theorem~\ref{thm:HKL}(b) for the pair $\{u,v\}$ (in whichever order makes
the lengths nonincreasing). First,
$u+v>\tfrac32n\ge\tfrac32\md(z)\ge\md(z)+\cy(z)$, so
$u+v>\md(z)+\cy(z)$. Second, $u+v$ is odd (one length even, one odd) and
$\md(z)+\cy(z)\equiv \md(z)-\cy(z)=\delta \pmod 2$ is
odd as well since $z$ is odd; so
$u+v\ge \md(z)+\cy(z)$ with the correct parity. Third, since
$u,v\in(\tfrac34n,\,n]$,
\[
|u-v|\;<\;\tfrac14 n\;\le\;\delta\;=\;\md(z)-\cy(z) .
\]
Both lengths are at least $2$ and at most $n$, and $z\neq1$ as $z$ is odd. Hence
$z=AB$ with $A$ a $u$-cycle and $B$ a $v$-cycle. Since $\gcd(u,r)=1$ and
$\gcd(v,s)=1$, Lemma~\ref{lem:root} gives
$x^r=A$, $y^s=B$, so $z=x^ry^s$.
\end{proof}

\section{Odd targets of small support: obstruction and escape}\label{sec:sparse}

This case is the heart of the paper: it is here that the two-single-cycle
method provably demands twin-prime-type input
(Proposition~\ref{prop:obstruction} below), and here that the ternary
construction enters.
Throughout this section we assume \eqref{eq:standing} and consider odd targets
$z$ with
\[
\md(z)<n/2,\qquad \delta(z)\ge 3 .
\]

\subsection{The two-cycle obstruction}\label{ssec:obstruction}
We first make precise the claim of the introduction that no analogue of
Propositions~\ref{prop:even} and \ref{prop:dense} --- two single cycles, one an
$r$-th power and one an $s$-th power --- can cover the present case. This is the
one place where the necessity direction of Theorem~\ref{thm:HKL} is used; nothing
later depends on it.

\begin{proposition}[two-cycle obstruction]\label{prop:obstruction}
Let $r$ be even and $s$ odd, and suppose every prime $p\le x$ divides $r$ and
every odd prime $p\le x$ divides $s$, where $x\ge 5$. Let $z\in S_n$ be a
$4$-cycle, and suppose $z=AB$ where $A$ is an $l_1$-cycle that is an $r$-th
power in $S_n$ and $B$ is an $l_2$-cycle that is an $s$-th power
($l_1,l_2\ge2$). Then $l_1$ is odd with every prime factor exceeding $x$, $l_2$
is even with every odd prime factor exceeding $x$, and $|l_1-l_2|\le 3$. In
particular, if $n<x^2$, then $l_1$ is a prime and $l_2$ is $2^a$ or $2^aq$ with
$q>x$ prime.
\end{proposition}

\begin{proof}
An $l$-cycle is a $t$-th power in $S_n$ if and only if $\gcd(l,t)=1$: each cycle
of $w$ of length $\ell$ contributes $\gcd(\ell,t)$ cycles of length
$\ell/\gcd(\ell,t)$ to $w^t$, so a single $l$-cycle arises only from a single
$\ell$-cycle with $\ell=l$ and $\gcd(l,t)=1$. Hence $\gcd(l_1,r)=1$ and
$\gcd(l_2,s)=1$, giving the prime-factor claims. Since $r$ is even, $A$ is an
even permutation, so $l_1$ is odd; since $z$ is odd, $l_2$ is even. A $4$-cycle
has one nontrivial cycle, so case (a) of Theorem~\ref{thm:HKL} is impossible, and
the necessity of case (b) gives $|l_1-l_2|\le \md(z)-\cy(z)=3$. If $n<x^2$, a
length at most $n$ whose (odd) prime factors all exceed $x$ has at most one of
them.
\end{proof}

Taking $r=\prod_{p\le x}p$ and $s=\prod_{2<p\le x}p$, so that
$\log m=\theta(x)\asymp x$, a two-single-cycle proof of the Main Theorem would
thus produce, for every sufficiently large $x$, primes $u,q\in(x,n]$ with
$u-2^aq\in\{\pm1,\pm3\}$ (or $u$ within $3$ of a power of $2$), where
$n\asymp x$: a constellation problem of twin-prime type, involving the dilate
$2^aq$ rather than a shift, and beyond any current method. (Of course the
$4$-cycle itself is an $s$-th power for odd $s$ --- the obstruction concerns the
two-single-cycle method with both cycles nontrivial, not membership in the image
of the word map.)

\subsection{The construction, given a good prime}
Call a prime $\tilde q$ \emph{good} (for the fixed data $r,s,n$) if
\begin{enumerate}
\item[(G1)] $n/4<\tilde q\le n/2$ and $\tilde q\nmid s$;
\item[(G2)] the odd number $M=2\tilde q-1$ admits a representation
$M=q_1+q_2+q_3$ with $q_1,q_2,q_3$ primes none of which divides $rs$.
\end{enumerate}

\begin{lemma}\label{lem:construction}
If a good prime $\tilde q$ exists, then every odd $z\in S_n$ with $\md(z)<n/2$ and
$\delta(z)\ge 3$ satisfies $z=x^ry^s$.
\end{lemma}

\begin{proof}
Put $N=2\tilde q\in(n/2,\,n]$; since $\md(z)<n/2<N\le n$, we may choose
$\Omega\subseteq\{1,\dots,n\}$ with $\supp(z)\subseteq\Omega$ and $|\Omega|=N$.
Let $q_1+q_2+q_3=N-1$ be as in (G2) and set
\[
\lambda=(q_1,q_2,q_3,1)\ \vdash\ N,\qquad \ell(\lambda)=4 .
\]
Let $\mu\vdash N$ be the type of $z$ on $\Omega$; it consists of the $\cy(z)$
nontrivial cycles of $z$ together with $N-\md(z)$ fixed points, so
$\ell(\mu)=\cy(z)+N-\md(z)=N-\delta(z)$. Then
\[
\ell(\lambda)+\ell(\mu)\;=\;N+4-\delta(z)\;\le\;N+1
\iff \delta(z)\ge 3,
\]
which holds; and modulo $2$,
\[
\ell(\lambda)+\ell(\mu)\;=\;N+4-\delta(z)\;\equiv\;N+1
\iff \delta(z)\equiv 3\equiv 1 ,
\]
which holds since $z$ is odd (Lemma~\ref{lem:elem}(1)). By
Theorem~\ref{thm:Boccara} there are $A\in C_\lambda$ and an $N$-cycle $B$ on
$\Omega$ with $AB=z$ (using, as noted after Theorem~\ref{thm:HKL}, the freedom to put the
$C_\lambda$-factor first). The nontrivial cycles of $A$ have prime lengths
$q_1,q_2,q_3$, none dividing $r$, so Lemma~\ref{lem:root} gives $x$ with $x^r=A$.
The factor $B$ is an $N$-cycle with $N=2\tilde q$, and $\gcd(N,s)=1$ since $s$ is
odd and $\tilde q\nmid s$; Lemma~\ref{lem:root} gives $y$ with $y^s=B$. Extending
by the identity off $\Omega$, $x^ry^s=AB=z$.
\end{proof}

\subsection{Averaging the singular series}
It remains to produce a good prime. For any \emph{single} candidate $M$ the sieve
bound of Theorem~\ref{thm:sieve} loses a factor $\asymp\log n$ against $r_3(M)$
when $rs$ has many prime divisors (see Remark~\ref{rem:avg}); the remedy is to
average over a full window of candidates. Let
\[
W\;=\;\{\tilde q \text{ prime}:\ (n+2)/4<\tilde q\le n/2\} ,
\]
the lower endpoint chosen so that $M=2\tilde q-1>n/2$ for every $\tilde q\in W$.

\begin{lemma}\label{lem:Wsize}
For $n\ge 10^8$, $\ |W|\ \ge\ 0.15\,n/\log n$.
\end{lemma}

\begin{proof}
By \eqref{eq:pilow}, $\pi(n/2)>(n/2)/\log(n/2)>n/(2\log n)$. By \eqref{eq:piup}
and $\log(n/4)\ge 0.9\log n$ (valid for $n\ge 4^{10}$),
$\pi((n+2)/4)\le\pi(n/4)+1<1.25506\,(n/4)/(0.9\log n)+1<0.349\,n/\log n$.
Subtract.
\end{proof}

\begin{lemma}[Singular series on average]\label{lem:avg}
There is an absolute constant $C_3$ (one may take $C_3=4$ for $n\ge 10^8$) such
that for every integer $b$ with $0\le b\le n+2$,
\[
\sum_{\substack{\tilde q\in W\\ 2\tilde q-b\ \mathrm{even},\ \ge 4}}
S(2\tilde q-b)\;\le\;C_3\,\frac{n}{\log n},
\]
with $S(\cdot)$ as in Theorem~\ref{thm:sieve}.
\end{lemma}

\begin{proof}
For even $k\ge 4$ expand the product: writing $g$ for the multiplicative function
supported on odd squarefree integers with $g(p)=1/(p-2)$,
\[
S(k)\;=\;\prod_{p\mid k,\,p>2}\Bigl(1+\frac1{p-2}\Bigr)
\;=\;\sum_{\substack{d\mid k\\ d \text{ odd squarefree}}} g(d) .
\]
Write $\Sigma$ for the sum in the statement; its arguments $k=2\tilde q-b$ all
satisfy $4\le k\le n$, so every divisor of such a $k$ is at most $n$. Hence,
interchanging sums,
\[
\Sigma
\;\le\;\sum_{\substack{d\le n\\ d \text{ odd squarefree}}} g(d)\,
\#\{\tilde q\in W:\ 2\tilde q\equiv b \pmod d\}.
\]
For each odd $d$ the congruence $2\tilde q\equiv b\pmod d$ determines a single
residue class $a_d$ modulo $d$. Split at $D=n^{1/3}$.

If $d\le D$: when $\gcd(a_d,d)>1$, at most $\omega(d)\le\log d$ primes lie in the
class; when $\gcd(a_d,d)=1$, Brun--Titchmarsh \eqref{eq:BT} with $x=n/2$ gives,
for $n\ge 10^8$,
\[
\#\{\cdot\}\;\le\;\pi(n/2;d,a_d)\;<\;\frac{n}{\varphi(d)\log(n/(2d))}
\;\le\;\frac{n}{0.6\,\varphi(d)\log n}
\]
using $\log(n/(2d))\ge\tfrac23\log n-\log2\ge0.6\log n$. Since
$\sum_{d \text{ odd squarefree}} g(d)/\varphi(d)
=\prod_{p>2}\bigl(1+\tfrac1{(p-2)(p-1)}\bigr)<1.75$,
the total contribution of $d\le D$ is at most
$\tfrac{1.75}{0.6}\,n/\log n+\sum_{d\le D}g(d)\log d
\le 2.92\,n/\log n+e^{1.2}\log^2\!D$,
using $\sum_{d\le x}g(d)\le\prod_{p\le x,\,p>2}(1+\tfrac1{p-2})\le e^{1.2}\log x$
for $x\ge 10$ (take logarithms and use
$\sum_{2<p\le x}1/(p-2)\le\log\log x+1.2$), so that
$\sum_{d\le D}g(d)\log d\le \log D\sum_{d\le D}g(d)\le e^{1.2}\log^2\!D
\le0.37\log^2 n$.

If $d>D$: crudely $\#\{\tilde q\in W:\ 2\tilde q\equiv b\pmod d\}\le n/(4d)+1$.
Since $\sum_d g(d)/\sqrt d=\prod_{p>2}\bigl(1+\tfrac1{\sqrt p\,(p-2)}\bigr)<2.33$,
\[
\sum_{D<d\le n}g(d)\Bigl(\frac n{4d}+1\Bigr)
\;\le\;\frac n4\cdot\frac{2.33}{\sqrt D}\;+\;\sum_{d\le n}g(d)
\;\le\;0.59\,n^{5/6}\;+\;e^{1.2}\log n .
\]
The three secondary terms $0.37\log^2n$, $0.59\,n^{5/6}$ and $e^{1.2}\log n$,
divided by $n/\log n$, are each decreasing for $n\ge10^8$, and at $n=10^8$ they
total less than $0.51\,n/\log n$; hence
$\Sigma\le 2.92\,n/\log n+0.51\,n/\log n<4\,n/\log n$.
\end{proof}

\begin{proposition}\label{prop:good}
There are absolute constants $C^*$ and $N^*$ (with $C^*$ effective given $C_2$ of
Theorem~\ref{thm:sieve} and $c_1$ of Theorem~\ref{thm:vino}, and $N^*$
ineffective) such that: if $n\ge N^*$ and $n\ge C^*L$, then a good prime
$\tilde q\in W$ exists.
\end{proposition}

\begin{proof}
We discard the bad elements of $W$ and count what remains. Set
$\tilde\omega=\#\{p\le n: p\mid rs\}\le\omega(m)$. Assume $n\ge C^*L$ with
$C^*\ge 100$ and $n\ge (C^*)^2$, so that $\log(n/C^*)\ge\tfrac12\log n$; by
Lemma~\ref{lem:omega},
\begin{equation}\label{eq:omegabound}
\tilde\omega\;\le\;1.28\,\frac{L}{\log L}+1429
\;\le\;\frac{2.56}{C^*}\cdot\frac{n}{\log n}+1429 .
\end{equation}
(For $L<e$ we have $m\le 15$, so $\omega(m)\le2$ and \eqref{eq:omegabound} is
trivial; for $e\le L\le\sqrt n$, monotonicity of $t/\log t$ on $[e,\infty)$
bounds the first term by $2.56\sqrt n/\log n$, which is even smaller. Thus
\eqref{eq:omegabound} holds in all cases for $n\ge (C^*)^2$.)

\emph{Bad set 1: $\tilde q\mid s$.} At most $\tilde\omega$ elements of $W$.

\emph{Bad set 2: primes $\tilde q$ for which many representations meet $rs$.}
For $\tilde q\in W$ put $M=2\tilde q-1\in(n/2,\,n]$, an odd number (the range by
the choice of $W$). Let $\mathrm{Bad}(M)$ be the number of
ordered representations $M=q_1+q_2+q_3$ in primes with some $q_i\mid rs$. Then
\[
\mathrm{Bad}(M)\;\le\;3\sum_{\substack{p\mid rs\\ p\le M}} r_2(M-p).
\]
The term $p=2$ contributes $r_2$ of an odd number, at most $2$, so at most $6$ per
$\tilde q$. For odd $p$ the argument $k=M-p$ is even; when $k<\sqrt n$ we use
$r_2(k)\le k\le\sqrt n$, and when $k\ge\sqrt n$ Theorem~\ref{thm:sieve} with
$k\le n$ and $\log^2k\ge\tfrac14\log^2n$ gives $r_2(k)\le 4C_2\,S(k)\,n/\log^2n$.
Summing over $\tilde q\in W$ and using Lemma~\ref{lem:avg} for each fixed $p$
(with $b=1+p$):
\[
\sum_{\tilde q\in W}\mathrm{Bad}(2\tilde q-1)
\;\le\;6|W|\;+\;3\,\tilde\omega\,\Bigl(\tfrac{\sqrt n}2+1\Bigr)\sqrt n
\;+\;12\,C_2\,\frac{n}{\log^2 n}\cdot \tilde\omega\cdot C_3\frac{n}{\log n} .
\]
By Theorem~\ref{thm:vino}, for $n$ large every odd $M\in(n/2,n]$ has
\[
r_3(M)\;\ge\;c_1M^2/\log^3M\;\ge\;\tfrac{c_1}4\,n^2/\log^3 n\;=:\;T .
\]
Let
$E=\{\tilde q\in W:\ \mathrm{Bad}(2\tilde q-1)\ge\tfrac12 r_3(2\tilde q-1)\}$;
then $|E|\le 2T^{-1}\sum_{\tilde q}\mathrm{Bad}$, i.e.
\[
|E|\;\le\;\frac{8\log^3n}{c_1n^2}\Bigl[6|W|+2\tilde\omega n
+12\,C_2C_3\,\tilde\omega\,\frac{n^2}{\log^3 n}\Bigr]
\;=\;\frac{96\,C_2C_3}{c_1}\,\tilde\omega\;+\;o\Bigl(\frac n{\log n}\Bigr),
\]
the secondary terms being $O(|W|\log^3 n/n^2)+O(\tilde\omega\log^3n/n)
=O(\log^2 n)$. By \eqref{eq:omegabound},
\[
\frac{96\,C_2C_3}{c_1}\,\tilde\omega
\;\le\;\frac{246\,C_2C_3}{c_1\,C^*}\cdot\frac n{\log n}
\;+\;\frac{96\cdot 1429\,C_2C_3}{c_1}.
\]
Now choose $C^*=\max\{100,\ 4920\,C_2C_3/c_1\}$, so the first term is at most
$0.05\,n/\log n$, and let $N^*$ be large enough that: $n\ge\max\{10^8,(C^*)^2\}$;
$2\tilde q-1\ge M_0$ for all $\tilde q\in W$; the constant $1429$ of
\eqref{eq:omegabound} is at most $0.004\,n/\log n$, so that
$\tilde\omega\le0.03\,n/\log n$; and each of the three remaining terms in the
bound for $|E|$ --- the constant $96\cdot1429\,C_2C_3/c_1$ and the two secondary
terms --- is at most $0.01\,n/\log n$. Then, using Lemma~\ref{lem:Wsize},
\[
\#\{\text{good }\tilde q\}\;\ge\;|W|-\tilde\omega-|E|
\;\ge\;\Bigl(0.15-0.03-0.05-0.04\Bigr)\frac n{\log n}\;>\;0
\]
(the $0.03$ covering Bad set~1, the $0.05$ the main term of $|E|$, and the
$0.04$ its three remaining terms).
Any $\tilde q\in W$ outside both bad sets satisfies (G1), and (G2) because
$r_3(2\tilde q-1)-\mathrm{Bad}(2\tilde q-1)\ge\tfrac12r_3(2\tilde q-1)>0$: some
representation avoids the primes of $rs$ entirely.
\end{proof}

\begin{remark}\label{rem:avg}
The averaging over $\tilde q$ is not cosmetic. For a single fixed $M$ the
worst-case bound $S(M-p)\ll\log\log n$, multiplied by up to $L/\log 2\asymp n/C$
possible $p$, exceeds $r_3(M)$ by a factor $\asymp\log n\log\log n/C$ --- and
even the sharpened count of Lemma~\ref{lem:omega} (attained when $rs$ is a
primorial) leaves a factor $\asymp\log\log n/C$: no fixed choice of $M$ can be
proved to work by these estimates alone. Averaging restores the loss
because, as $\tilde q$ varies over a full window of primes, the singular series
$S(2\tilde q-1-p)$ is bounded on average (Lemma~\ref{lem:avg}); and the sharpened
count $\tilde\omega\ll L/\log L$ of Lemma~\ref{lem:omega} --- distinct prime
divisors cannot all be small --- supplies exactly the factor $\log n$ that the
trivial bound $\tilde\omega\le L/\log 2$ would lose. As stressed in the
introduction, the quantifier order is what makes the averaging legitimate: $r$
and $s$ are given before $\tilde q$ is chosen.
\end{remark}

\section{Proof of the Main Theorem}\label{sec:proof}

\begin{proof}[Proof of the Main Theorem]
Let $C=\max\{4.65,\ C^*\}$ and $N_0=\max\{381700,\ N^*\}$ (the maxima record what
the effective cases alone require) with $C^*,N^*$ as in
Proposition~\ref{prop:good}, and suppose $n\ge\max\{N_0,\,C\log m\}$. By the
reductions of \S\ref{sec:red} we may assume \eqref{eq:standing}; the targets
$z=1$ and $z$ a transposition are done there. Let $z\in S_n$ be any other target.

If $z$ is even, Proposition~\ref{prop:even} applies. If $z$ is odd, then
$\delta(z)$ is odd; since $z$ is not a transposition, $\delta(z)\ge3$
(Lemma~\ref{lem:elem}(3)). If $\md(z)\ge n/2$, Proposition~\ref{prop:dense}
applies. Otherwise $\md(z)<n/2$, and Proposition~\ref{prop:good} together with
Lemma~\ref{lem:construction} applies. In every case $z=x^ry^s$ for some
$x,y\in S_n$.
\end{proof}

\section{Optimality of the logarithmic shape}\label{sec:lower}

Brenner, Evans and Silberger showed that on $A_n$ the constant of their theorem
cannot be taken to be $1$: for $r=s=n!$ every value of $x^ry^s$ on $A_n$ is
trivial, while $\theta(n)=\log m<n$ for infinitely many $n$ \cite[proof of
Theorem~3$'$]{BES}. Both exponents there are even, so this says nothing about
the symmetric-group problem, in which one exponent must stay odd. The following
variant respects the parity hypothesis.

\begin{proposition}\label{prop:lower}
Let $n\ge3$, let $r=\operatorname{lcm}(1,2,\dots,n)$ and let $s$ be the odd part
of $r$. Then $r$ is even, $s$ is odd, $\log m(r,s)=\theta(n)$, and every value
of $x^ry^s$ on $S_n$ has all cycle lengths powers of $2$. In particular no
$3$-cycle is a value, and $x^ry^s$ is not universal on $S_n$.
\end{proposition}

\begin{proof}
Every cycle length $\ell\le n$ divides $r$, so by the power-of-a-cycle fact
quoted before Lemma~\ref{lem:root} every cycle of every $x\in S_n$ collapses to
fixed points: $x^r=1$ identically, and the values of the word are exactly the
$s$-th powers. Write a cycle length of $y$ as $\ell=2^ab$ with $b$ odd. Then
$b\le n$ gives $b\mid r$, and $b$ odd gives $b\mid s$ ($s$ is the odd part of
$r$), while $s$ odd gives $\gcd(\ell,s)=b$; so, by the
same fact, the cycle contributes $b$ cycles of length $2^a$ to $y^s$. Thus
every cycle length of $y^s$ is a power of $2$, and a $3$-cycle ($n\ge3$) is not
an $s$-th power. Finally, the primes dividing $rs$ are exactly the primes up to
$n$, so $\log m(r,s)=\theta(n)$.
\end{proof}

\begin{corollary}\label{cor:lower}
There is no $N_0$ such that, for all nonzero integers $r,s$ with at least one
odd, the word $x^ry^s$ is universal on $S_n$ whenever
$n\ge\max\{N_0,\,\log m(r,s)\}$. Thus the constant $C$ of the Main Theorem
cannot be taken to be $1$ (nor, a fortiori, smaller), and the logarithmic shape
is optimal for $S_n$ just as for $A_n$.
\end{corollary}

\begin{proof}
There are infinitely many integers $n$ with $\theta(n)<n$ (see \cite[p.~67]{RS},
as used in the proof of \cite[Theorem~3$'$]{BES}). Given a candidate $N_0$, pick
such an $n\ge\max\{N_0,3\}$ and take $r,s$ as in
Proposition~\ref{prop:lower}: then $n\ge\max\{N_0,\log m(r,s)\}$, yet $x^ry^s$
is not universal on $S_n$.
\end{proof}

\section{Concluding remarks}\label{sec:remarks}

The two 1986 \emph{Proceedings} papers and Merzlyakov's problem set a precise
program: a degree bound for $x^ry^s$ logarithmic in $m(r,s)$, uniform in the
word, on each of $A_n$ and $S_n$, together with the matching lower bound. With
the Main Theorem and \S\ref{sec:lower}, both halves of the symmetric-group case
are now closed, forty years after the problem was posed.

Beyond the statement, we would single out the diagnosis of
\S\ref{ssec:obstruction} as a contribution in its own right:
Proposition~\ref{prop:obstruction} locates the difficulty of Problem~10.32 on
the divide between binary and ternary additive prime problems --- two cycle
lengths force a twin-prime-type constellation, out of reach, while a factor
shape with three prime summands moves the problem to ternary Goldbach, where
Vinogradov's theorem applies. The analytic devices themselves are classical ---
the three-primes lower bound, an upper-bound sieve, the averaging of singular
series over a window --- and we claim no novelty for any of them individually.
What may transfer is the arrangement: a three-primes theorem relative to an
excluded set of primes that is fixed before the window is chosen, with the
adversary defeated on average rather than pointwise, and the count
$\omega(m)\ll\log m/\log\log m$ closing the final factor. That pattern ---
choose the factor shape to fit the arithmetic, then let an exact factorization
criterion certify it --- may be useful, as may the dyadic windows of
\S\ref{sec:dense}, in other covering problems in symmetric groups where
prescribed prime-avoidance is the obstacle.


\begin{remark}[Shapes of the factors]
The proof gives more than universality: every $z\in S_n$ is expressed with factors
of an explicit and very restricted shape --- two prime-length cycles for even $z$;
an odd-prime cycle times a $(2^a\cdot\text{prime})$-cycle for odd $z$ of large
support; and a type-$(q_1,q_2,q_3,1)$ permutation times a full even cycle on an
$N$-point subset for odd $z$ of small support. The analysis in
\S\ref{sec:intro} explains why some change of shape in the last case is not an
artifact of the method: for odd targets with bounded $\delta(z)$ and adversarial
$r,s$, the Herzog--Kaplan--Lev constraints force the two cycle lengths within
bounded distance of each other while avoiding prescribed primes --- twin-prime
territory (Proposition~\ref{prop:obstruction}). It is the passage from two
factors' worth of freedom to
three summands' worth --- from binary to ternary Goldbach --- that makes the
small-support case provable.
\end{remark}

\begin{remark}[Effectivity]\label{rem:effective}
The constant $C$ is an explicit number once the sieve constant $C_2$ and
Vinogradov constant $c_1$ are fixed from the cited sources; no attempt has been
made to optimize it. The threshold $N_0$ is ineffective solely through $M_0$ of
Theorem~\ref{thm:vino}. Replacing Theorem~\ref{thm:vino} by the explicit
three-primes estimates underlying Helfgott's proof \cite{Helfgott} would, at the
cost of numerical work orthogonal to the ideas here, render $N_0$ effective as
well. By contrast, targets that are even or move at least $n/2$ points need only
Propositions~\ref{prop:even} and \ref{prop:dense}: for these,
$n\ge\max\{381700,\,4.65\log m\}$ suffices, fully explicitly.
\end{remark}

\begin{remark}[Constants]
The case constants $4.5$ and $4.65$ come from windows of relative width
$\asymp1/4$, which is the natural floor of the present placement: sharper
Chebyshev bounds (Dusart's, for instance) push both constants toward $4$, but
not below it, without a structural change. Brenner, Evans and Silberger
\cite{BES} obtain every $c>8/5$ for $A_n$ by a finer argument in which whole
ranges of primes dividing $m$ are played against each other, and
\S\ref{sec:lower} shows $c>1$ is required for $S_n$ as well. Locating the
optimal constant --- for the effective cases somewhere in $(1,\,4.65]$, and for
the theorem as a whole tied to the Vinogradov bookkeeping of \S\ref{sec:sparse}
--- remains open, as does an $S_n$ analogue of the sharp $8/5$ threshold of
\cite{BES}.
\end{remark}

\begin{remark}[Verification]\label{rem:verification}
Beyond Remark~\ref{rem:machine}, two computations were run against this paper's
own argument, chiefly to guard against misquotation of the imported criteria and
against slips in composition-order conventions: the condition chains fed to
Theorems~\ref{thm:HKL} and \ref{thm:Boccara} in \S\S4--6 were checked on $600$
random instances $(n,r,s,z)$ with $n\le3000$, and the small-support construction
of Lemma~\ref{lem:construction} was executed end to end in $S_{14}$--$S_{16}$ on
$400$ random targets, the identity $x^ry^s=z$ holding exactly in every instance.
The lower-bound construction of \S\ref{sec:lower} was likewise verified
exhaustively for $n\le7$. The scripts, including a rigorous certification of the
numerical constants of Lemma~\ref{lem:avg}, are available at
\url{https://github.com/AviKndr/universal-words} (\texttt{validate.py}).
\end{remark}

\section*{Acknowledgements}
This paper was researched and written with substantial assistance from Claude
(Anthropic), which was used for literature retrieval, for the design and checking
of the argument, and for drafting. The two factorization criteria and the
constructions of \S\S4--6 were additionally checked by machine computation as
indicated in Remarks~\ref{rem:machine} and \ref{rem:verification}. Errors are the
author's. The statement of Problem~10.32 is quoted from the Kourovka Notebook
\cite{KN}.

\begin{thebibliography}{99}

\bibitem{Bertram} E.~Bertram, \emph{Even permutations as a product of two
conjugate cycles}, J. Combin. Theory Ser.~A \textbf{12} (1972), 368--380.

\bibitem{Boccara80} G.~Boccara, \emph{Nombre de repr\'esentations d'une
permutation comme produit de deux cycles de longueurs donn\'ees}, Discrete Math.
\textbf{29} (1980), 105--134.

\bibitem{Boccara82} G.~Boccara, \emph{Cycles comme produit de deux permutations
de classes donn\'ees}, Discrete Math. \textbf{38} (1982), 129--142.

\bibitem{BES} J.~L.~Brenner, R.~J.~Evans and D.~M.~Silberger, \emph{The
universality of words $x^ry^s$ in alternating groups}, Proc. Amer. Math. Soc.
\textbf{96} (1986), no.~1, 23--28.

\bibitem{Droste} M.~Droste, \emph{On the universality of words for the
alternating groups}, Proc. Amer. Math. Soc. \textbf{96} (1986), no.~1, 18--22.

\bibitem{EKS} A.~L.~Edmonds, R.~S.~Kulkarni and R.~E.~Stong, \emph{Realizability
of branched coverings of surfaces}, Trans. Amer. Math. Soc. \textbf{282} (1984),
773--790.

\bibitem{HR} H.~Halberstam and H.-E.~Richert, \emph{Sieve Methods}, Academic
Press, 1974.

\bibitem{Helfgott} H.~A.~Helfgott, \emph{The ternary Goldbach conjecture is
true}, arXiv:1312.7748 (2013); expanded as \emph{The ternary Goldbach problem},
arXiv:1501.05438.

\bibitem{HKL} M.~Herzog, G.~Kaplan and A.~Lev, \emph{Representation of
permutations as products of two cycles}, Discrete Math. \textbf{285} (2004),
323--327.

\bibitem{KN} \emph{Unsolved problems in group theory. The Kourovka Notebook},
21st ed., E.~I.~Khukhro and V.~D.~Mazurov (eds.), Sobolev Institute of
Mathematics, Novosibirsk, 2026; arXiv:1401.0300. Problem~10.32.

\bibitem{LaSh} M.~Larsen and A.~Shalev, \emph{Word maps and Waring type
problems}, J. Amer. Math. Soc. \textbf{22} (2009), 437--466.

\bibitem{LST} M.~Larsen, A.~Shalev and P.~H.~Tiep, \emph{The Waring problem
for finite simple groups}, Ann. of Math. (2) \textbf{174} (2011), no.~3,
1885--1950.

\bibitem{LOST} M.~W.~Liebeck, E.~A.~O'Brien, A.~Shalev and P.~H.~Tiep,
\emph{The Ore conjecture}, J. Eur. Math. Soc. \textbf{12} (2010), no.~4,
939--1008.

\bibitem{MV} H.~L.~Montgomery and R.~C.~Vaughan, \emph{The large sieve},
Mathematika \textbf{20} (1973), 119--134.

\bibitem{Petronio} C.~Petronio, \emph{The Hurwitz existence problem for surface
branched covers}, Winter Braids Lect. Notes \textbf{7} (2020), Course~II, 1--43.

\bibitem{RS} J.~B.~Rosser and L.~Schoenfeld, \emph{Approximate formulas for some
functions of prime numbers}, Illinois J. Math. \textbf{6} (1962), 64--94.

\bibitem{RS75} J.~B.~Rosser and L.~Schoenfeld, \emph{Sharper bounds for the
Chebyshev functions $\theta(x)$ and $\psi(x)$}, Math. Comp. \textbf{29} (1975),
243--269.

\bibitem{Thom} R.~Thom, \emph{L'\'equivalence d'une fonction diff\'erentiable et
d'un polyn\^ome}, Topology \textbf{3} (1965), suppl.~2, 297--307.

\bibitem{Vaughan} R.~C.~Vaughan, \emph{The Hardy--Littlewood Method}, 2nd ed.,
Cambridge Tracts in Math. \textbf{125}, Cambridge Univ. Press, 1997.

\bibitem{Vino} I.~M.~Vinogradov, \emph{Representation of an odd number as a sum
of three primes}, C. R. (Dokl.) Acad. Sci. URSS \textbf{15} (1937), 169--172.

\bibitem{Wilf} H.~S.~Wilf, \emph{generatingfunctionology}, 2nd ed., Academic
Press, 1994.

\end{thebibliography}

\end{document}
